\input{configTD}


\title{Probabilités TD3 - Modélisations avec des lois continues}

\author{}
\date{}

\newenvironment{solution}
    {\par\vspace{0.5em}\begin{mdframed}[linewidth=0.5pt]\par}
    {\end{mdframed}\par\vspace{0.5em}}

\begin{document}
\sffamily
\maketitle

\section*{1 - Espérance et variance}
\begin{enumerate}
\item On rappelle la formule de l'espérance pour une variable continue : 
\[
E(X) = \int_{-\infty}^\infty x f_X(x) \, dx\,.
\]
Cette intégrale fait la moyenne pondérée des valeurs prises par $X$, avec les probabilités associées.\newline 
Exrpimez l'espérance de $X^2$ par analogie avec l'intégrale ci-dessus.
\ifthenelse{\boolean{showSolutions}}{
    \begin{solution}
        $$
        E(X^2) = \int_{-\infty}^\infty x^2 f_X(x) \, dx
        $$
    \end{solution}
}{}
\item Rappeler la densité de la loi uniforme sur l'intervalle $[-5, 5]$ puis calculer la variance d'une variable aléatoire $Y$ suivant cette loi. 
\ifthenelse{\boolean{showSolutions}}{
    \begin{solution}
        La densité de la loi uniforme sur l'intervalle $[-5, 5]$ est :
        $$
        f_Y(y) = \frac{1}{10} \mathbf{1}_{[-5, 5]}(y)
        $$

        Calculons la variance :
        \begin{align*}
            \text{Var}(Y) &= E\big(\,Y^2\, \big) - E\big(\,Y\, \big)^2 \\
        \end{align*}
        On peut calculer $E\big(\,Y^2\, \big)$ en intégrant $y^2$ sur l'intervalle $[-5, 5]$ et en multipliant par la densité de la loi uniforme :
        $$
        E\big(\,Y^2\, \big) = \frac{1}{10} \int_{-5}^5 y^2 \, dy = \frac{1}{10} \left[ \frac{y^3}{3} \right]_{-5}^5 = \frac{1}{10} \left( 2\,  \frac{5^3}{3} \right) = \frac{25}{3}\,.
        $$
        On peut calculer $E\big(\,Y\, \big)$ en intégrant $y$ sur l'intervalle $[-5, 5]$ et en multipliant par la densité de la loi uniforme :
        $$
        E\big(\,Y\, \big) = \frac{1}{10} \int_{-5}^5 y \, dy = \frac{1}{10} \left[ \frac{y^2}{2} \right]_{-5}^5 = \frac{1}{10} \left( \frac{25}{2} - \frac{25}{2} \right) = 0
        $$
        Donc :
        $$
        \text{Var}(Y) = \frac{25}{3} - 0^2 = \frac{25}{3}
        $$
    \end{solution}
}{}
\item Rappelez la densité, puis calculez la variance d'une variable aléatoire $Z$ suivant une loi exponentielle de paramètre~$\lambda$. 
\ifthenelse{\boolean{showSolutions}}{ 
    \begin{solution}
        La densité de la loi exponentielle de paramètre $\lambda$ est :
        $$
        f_Z(z) = \lambda e^{-\lambda z} \mathbf{1}_{z>0}
        $$
        Calculons la variance :
        \begin{align*}
            \text{Var}(Z) &= E\big(\,Z^2\, \big) - E\big(\,Z\, \big)^2 \\
        \end{align*}
        On peut calculer $E\big(\,Z^2\, \big)$ en intégrant $z^2$ sur l'intervalle $[0, \infty)$ et en multipliant par la densité de la loi exponentielle :
        $$
        E\big(\,Z^2\, \big) = \int_0^\infty z^2 \lambda e^{-\lambda z} \, dz
        $$
        Deux intégrations par parties successives permettent de calculer cette intégrale :
        \begin{align*}
            E\big(\,Z^2\, \big) &= \left[ -\frac{z^2}{\lambda} e^{-\lambda z} \right]_0^\infty + \int_0^\infty \frac{2z}{\lambda} e^{-\lambda z} \, dz \\
            &= 0 + \int_0^\infty \frac{2z}{\lambda} e^{-\lambda z} \, dz \\
            &= \left[ -\frac{2z}{\lambda^2} e^{-\lambda z} \right]_0^\infty + \int_0^\infty \frac{2}{\lambda^2} e^{-\lambda z} \, dz \\
            &= 0 + \frac{2}{\lambda^2} \int_0^\infty e^{-\lambda z} \, dz \\
            &= \frac{2}{\lambda^2}
        \end{align*}

        Dans le cours, on a calculé $E\big(\,Z\, \big) = \frac{1}{\lambda}$.

        Finalement, on a :
        $$
        \text{Var}(Z) = \frac{1}{\lambda^2}
        $$
    \end{solution}
}{}
\end{enumerate}
\vspace{1em}

\section*{2 - Fonctions de variables aléatoires - cas continu}

Soit la fonction $f$ définie sur $\mathbf{R}$ par : 
$$
f(x)=x e^{-\frac{x^2}{2}} \mathbf{1}_{x>0}
$$
\begin{enumerate}
    \item Montrez que $f$ est une densité de probabilité.
    \ifthenelse{\boolean{showSolutions}}{
        \begin{solution}
            On calcule l'intégrale :
            $$
            \int_0^\infty x e^{-\frac{x^2}{2}} \, dx 
            $$
            On connait une primitive de cette fonction : 
            \[
            \int_0^\infty x e^{-\frac{x^2}{2}} \, dx = \left[ -e^{-\frac{x^2}{2}} \right]_0^\infty = 1
            \]
            De plus, $f$ est positive et continue sur $\mathbf{R}$ donc $f$ est une densité de probabilité.
        \end{solution}
    }{}
    \vspace{1em}


    \item Soit $X$ une variable aléatoire continue dont la loi a pour densité $f$ et $Y=X^2$.\newline
    Exprimer $P\big(\,Y \in[a, b]\,\big)$ en fonction de la probabilité que $X$ se situe dans un certain intervalle.
    \ifthenelse{\boolean{showSolutions}}{
        \begin{solution}
            On a :
            $$
            P(Y \in [a, b]) = P(X^2 \in [a, b]) = P( X \in [\sqrt{a}, \sqrt{b}])
            $$
        \end{solution}
    }{}

    \item En déduire que
    \[
    P(Y \in [a, b]) = \int_{a}^{b} \frac{1}{2}e^{-\frac{x}{2}} \, dx
    \]
    \ifthenelse{\boolean{showSolutions}}{
        \begin{solution}
            On a :
            \begin{align*}
                P(Y \in [a, b]) 
                &= P( X \in [\sqrt{a}, \sqrt{b}]) \\
                &= \int_{\sqrt{a}}^{\sqrt{b}} xe^{-\frac{x^2}{2}} \, dx
            \end{align*}
            On effectue le changement de variable $u = x^2$, on obtient :
            $$
            P(Y \in [a, b]) = \int_{a}^{b} \frac{1}{2}e^{-\frac{u}{2}} \, du
            $$
        \end{solution}
    }{}
    \item En déduire la loi de $Y$, puis rappeler l'espérance et la variance de $Y$.
    \ifthenelse{\boolean{showSolutions}}{
        \begin{solution}
            La loi de $Y$ est une loi exponentielle de paramètre $1/2$. On rappelle que :
            $$
            E(Y) = \frac{1}{\lambda} = 2
            $$
            $$
            \text{Var}(Y) = \frac{1}{\lambda^2} = 4
            $$
        \end{solution}
    }{}
\end{enumerate}
\vspace{1em}



\section*{3 - Fonction de répartition}
La fonction de répartition d'une variable aléatoire $X$ est définie par :
$$
F_X(x) = P(X\leq x)\,.
$$

\begin{enumerate}
\item Exprimez $F_X$ à l'aide d'une intégrale et de la fonction de densité $f_X$.
\ifthenelse{\boolean{showSolutions}}{
    \begin{solution}
        Nous avons vu dans le cours que :
        \[
        P(X \in [a, b]) = \int_a^b f_X(t) \, dt
        \]
        Donc :
        $$
        F_X(x) = \int_{-\infty}^x f_X(t) \, dt
        $$
    \end{solution}
}{}
\item Exprimez la densité $f_X$ à l'aide de la dérivée de $F_X$.
\ifthenelse{\boolean{showSolutions}}{
    \begin{solution}
        $$
        f_X(x) = F_X'(x)
        $$
        La densité est la dérivée de la fonction de répartition.
    \end{solution}
}{}
\item Donner une égalité reliant $P(X > x)$ et $F_X(x)$.
\ifthenelse{\boolean{showSolutions}}{
    \begin{solution}
        $$
        P(X > x) = 1 - F_X(x)
        $$
    \end{solution}
}{}

\item Donner une égalité reliant $P(X \in [a, b])$ et $F_X(b) - F_X(a)$.
\ifthenelse{\boolean{showSolutions}}{
    \begin{solution}
        $$
        P(X \in [a, b]) = F_X(b) - F_X(a)
        $$
    \end{solution}
}{}
\item Calculez la fonction de répartition d'une variable aléatoire suivant une loi uniforme sur l'intervalle $[0,1]$.
\ifthenelse{\boolean{showSolutions}}{
    \begin{solution}
        Si $x$ est compris entre $0$ et $1$, on a :
        $$
        F_X(x) = \int_0^x 1 \, dt = x
        $$
        Si $x$ est supérieur à $1$, on a :
        $$
        F_X(x) = \int_0^1 1 \, dt = 1
        $$
        Si $x$ est inférieur à $0$, on a :
        $$
        F_X(x) = 0
        $$
    \end{solution}
}{}
\item Calculez la fonction de répartition d'une variable aléatoire suivant une loi exponentielle de paramètre $\lambda$.
\ifthenelse{\boolean{showSolutions}}{
    \begin{solution}
        On a, si $x$ est positif :
        $$
        F_X(x) = \int_0^x \lambda e^{-\lambda t} \, dt
        $$
        On connait une primitive de cette fonction :
        \[
        F_X(x) = \left[ -e^{-\lambda t} \right]_0^x = 1 - e^{-\lambda x}
        \]
        Cette fonction tend bien vers $1$ lorsque $x$ tend vers l'infini.
        Si $x$ est négatif, on a :
        $$
        F_X(x) = 0
        $$
        
    \end{solution}
}{}
\end{enumerate}

Dans certains cas, il sera plus pratique de travailler avec la fonction de répartition plutôt que la densité, comme dans l'exercice suivant par exemple :  
\newpage 
$\,$

\vspace{1em}

\section*{4 - Fonctions de répartitions, applications}

Soient $X_1, X_2 \sim \operatorname{Exp}(\lambda)$ des variables exponentielles indépendantes. 
On pose $M=\max \left(X_1, X_2\right)$.

\begin{enumerate}
\item Montrer que la fonction de répartition de $M$ en fonction de celle de $X_1$ et $X_2$ est :
$$
F_M(x) = (1 - e^{-\lambda x})^2
$$
\ifthenelse{\boolean{showSolutions}}{
    \begin{solution}
        On a :
        $$
        F_M(x) = P(M \leq x) = P(\max(X_1, X_2) \leq x) = P(X_1 \leq x\text{ et } X_2 \leq x)
        $$
        Or, $X_1$ et $X_2$ sont indépendantes, donc :
        $$
        F_M(x) = P(X_1 \leq x) P(X_2 \leq x) = (1 - e^{-\lambda x})(1 - e^{-\lambda x}) = (1 - e^{-\lambda x})^2
        $$
    \end{solution}
}{}
\item En déduire la densité de $M$.
\ifthenelse{\boolean{showSolutions}}{
    \begin{solution}
        $$
        f_M(x) = F_M'(x) = 2\lambda e^{-\lambda x}(1 - e^{-\lambda x})
        $$
    \end{solution}
}{}

On pose $m=\min \left(X_1, X_2\right)$.

\item Montrer que la fonction de répartition de $m$ en fonction de celle de $X_1$ et $X_2$ est :
$$
F_m(x) = 1 - e^{-2\lambda x}
$$
On pourra passer au complémentaire. 
\ifthenelse{\boolean{showSolutions}}{
    \begin{solution}
        On a :
        $$
        F_m(x) = P(m \leq x) = P(\min(X_1, X_2) \leq x) = P(X_1 \leq x\text{ ou } X_2 \leq x)
        $$
        En passant au complémentaire, on a :
        $$
        F_m(x) = 1 - P(X_1 > x \text{ et } X_2 > x)
        $$
        Or, $X_1$ et $X_2$ sont indépendantes, donc :
        $$
        F_m(x) = 1 - P(X_1 > x) P(X_2 > x) = 1 - (1 - F_{X_1}(x))(1 - F_{X_2}(x))
        $$
        On a $F_{X_1}(x) = 1 - e^{-\lambda x}$ et $F_{X_2}(x) = 1 - e^{-\lambda x}$, donc :
        $$
        F_m(x) = 1 - (1 - (1 - e^{-\lambda x}))(1 - (1 - e^{-\lambda x})) = 1 - (e^{-\lambda x})^2 = 1 - e^{-2\lambda x}
        $$
    \end{solution}
}{}    
\item En déduire la densité de $m$ et identifier sa loi. 
\ifthenelse{\boolean{showSolutions}}{
    \begin{solution}
        $$
        f_m(x) = F_m'(x) = 2\lambda e^{-2\lambda x}
        $$
        Cette densité est celle d'une loi exponentielle de paramètre $2\lambda$.
    \end{solution}
}{}


\vspace{2em}
On modélise deux files d'attente indépendantes. par des variables aléatoires $X_1$ et $X_2$ suivant une loi exponentielle de paramètre $1/10$ minutes.
\item Quelle est la probabilité que la file qui finit en premier soit la file 1 ?
\ifthenelse{\boolean{showSolutions}}{
    \begin{solution}
        Par symétrie, comme les deux files suivent la même loi et sont indépendantes, la probabilité que la file qui finit en premier soit la file 1 est $1/2$.
    \end{solution}
}{}
\item Quelle est la probabilité que la file qui finit en premier mette moins de 10 minutes ? 
\ifthenelse{\boolean{showSolutions}}{
    \begin{solution}
        $$
        P(m < 10) = F_m(10) = 1 - e^{-2\lambda \times 10} = 1 - e^{-2} \approx 86 \%\,.
        $$
    \end{solution}
}{}
\end{enumerate}

\vspace{1em}

\section*{5 - Temps d'attente à un guichet}
Dans une agence, on modélise le temps d'attente (en minutes) d'un client avant sa prise en charge par une variable aléatoire $T$ suivant une loi exponentielle de paramètre $\lambda > 0$.

On rappelle que la densité de $T$ est
$$
f_T(t)=\lambda e^{-\lambda t}\mathbf{1}_{t\geq 0}.
$$

L'agence considère qu'une attente supérieure à 10 minutes est pénalisante pour la satisfaction des clients.

\begin{enumerate}
    \item Vérifier que $f_T$ est bien une densité de probabilité.
    \item Déterminer la fonction de répartition de $T$.
    \item Quelle est la probabilité qu'un client attende plus de 10 minutes ?
    \item Quelle est la probabilité qu'un client attende entre 5 et 10 minutes ?
    \item Déterminer le temps d'attente moyen d'un client.
    \item On suppose qu'une attente supérieure à 10 minutes coûte en moyenne 8 euros à l'agence en geste commercial ou perte d'image, par client concerné. 
    Si l'ouverture d'un second guichet permet de doubler le paramètre et de passer de $\lambda$ à $2\lambda$, à quelle condition sur $\lambda$ devient-il rentable d'ouvrir ce second guichet, sachant que son coût est de 1 euro par client ?
\end{enumerate}

\end{document}